%% ============================================================
%% DocVerify — Pattern Recognition Letters (PRL) submission
%%
%% Format: elsarticle, 5p two-column, author-year (elsarticle-harv)
%% Style:  prletters.sty as required by the journal
%% Limit:  7 pages including references
%%
%% Build:
%%   pdflatex main.tex
%%   bibtex main
%%   pdflatex main.tex
%%   pdflatex main.tex
%% ============================================================

\documentclass[final,5p,times,twocolumn,number]{elsarticle}

%% ── Packages ─────────────────────────────────────────────────
\usepackage[T1]{fontenc}
\usepackage[utf8]{inputenc}
\usepackage{graphicx}
\usepackage{amsmath}
\usepackage{amssymb}
\usepackage{booktabs}
\usepackage{xcolor}
\usepackage{url}
\usepackage{microtype}

%% ── PRL journal style (required by the journal) ──────────────
\usepackage{prletters}
\usepackage{etoolbox}

%% Make abstract span full width (remove Article history / keywords panel)
\makeatletter
\renewcommand{\articleinfobox}{}
\patchcmd{\MaketitleBox}%
  {\parbox[t]{.4\textwidth}{\unvbox\absbox}}%
  {\parbox[t]{.777\textwidth}{\unvbox\absbox}}%
  {}{}
\patchcmd{\pprintMaketitle}%
  {\parbox[t]{.4\textwidth}{\unvbox\absbox}}%
  {\parbox[t]{.777\textwidth}{\unvbox\absbox}}%
  {}{}
\renewenvironment{abstract}{%
  \global\setbox\absbox=\vbox\bgroup%
  \hsize=.777\textwidth}%
 {\par\vspace{6pt}%
  \noindent\mbox{~}\hfill\copyright~\the\year\
  Elsevier~Ltd.~All~rights~reserved.\egroup}
\makeatother

%% ── hyperref ─────────────────────────────────────────────────
\usepackage[colorlinks=true, linkcolor=blue, citecolor=blue, urlcolor=blue]{hyperref}

%% ── Convenience macros ───────────────────────────────────────
\DeclareMathOperator*{\argmin}{arg\,min}
\newcommand{\TODO}[1]{\textcolor{red}{\textbf{[TODO: #1]}}}

%% Round-1 (first major revision) changes are now ACCEPTED: \rev renders as
%% normal text. Round-2 (this minor revision) changes are highlighted in blue
%% via \revB, the same colour used in the previous round, as the journal requires.
\newcommand{\rev}[1]{#1}
\newcommand{\revB}[1]{#1}
% Camera-ready: also make \revB plain -> \newcommand{\revB}[1]{#1}

%% Journal declaration
\journal{Pattern Recognition Letters}

% ============================================================
\begin{document}

\begin{frontmatter}

\title{Multi-objective Pareto hyperparameter optimization for joint detection and
       localization of forged identity documents}

\author[1]{Hern\'{a}n D\'{i}az Rodr\'{i}guez\corref{cor1}}
\cortext[cor1]{Corresponding author.}
\ead{diazhernan@uniovi.es}

\address[1]{Department of Computer Science, University of Oviedo, Asturias, Spain}

\begin{abstract}
Automated forensic analysis of identity documents requires simultaneously
detecting whether a document has been tampered with and identifying which
regions were altered. These two objectives are structurally coupled but exist
in tension: optimizing one with a fixed loss weight may degrade the other.
We propose DocVerify, a multi-task convolutional network paired with a
multi-objective hyperparameter optimization (HPO) framework that treats
detection quality (PR-AUC) and localization quality (Dice) as simultaneous
objectives. Using the Multiobjective Tree-structured Parzen Estimator (MOTPE)
within a 10-fold nested cross-validation protocol (50 trials per outer fold,
500 unique configurations evaluated across 5 inner splits, totaling 2\,500
individual model trainings), DocVerify selects configurations from the
Pareto-optimal front by minimizing Euclidean distance to the ideal point
(1, 1), without any manual weighting. Evaluated on the FantasyID benchmark,
DocVerify achieves PR-AUC = 0.9921 ± 0.0058 and
Dice = 0.856 ± 0.030 under nested cross-validation, and detection
F1@0.5 = 0.969 ± 0.014 with localization F1 = 0.807 ± 0.096 on the
blind test set (30 seeds). An ablation confirms that Pareto-optimized loss
weighting provides a statistically significant Dice improvement over
equal-weight joint training (p < 0.001, d = 1.01, large effect)
while maintaining equivalent detection accuracy. \revB{These figures characterize
in-domain performance on FantasyID; zero-shot transfer to the out-of-domain SIDTD
benchmark falls to near chance (PR-AUC 0.555), so the reported accuracy should be
read as bounded to the training domain.}
\end{abstract}

\begin{keyword}
  identity document forensics \sep
  multi-objective optimization \sep
  hyperparameter optimization \sep
  multi-task learning \sep
  document forgery detection
\end{keyword}

\end{frontmatter}

% ============================================================
% SECTIONS
% ============================================================

% ============================================================
% Section 1: Introduction
% ============================================================

\section{Introduction}
\label{sec:intro}

Identity document fraud poses a growing threat to national security, financial
integrity, and border control worldwide. The widespread availability of
image-editing software has made it possible to produce convincing forgeries of
passports, national identity cards, and driver's licenses at low cost.
Automated forensic analysis of identity documents
is therefore a critical component of modern digital identity verification
pipelines.

Detection of forged identity documents is inherently a two-level problem. At the
\emph{document level}, a system must decide whether a presented document is
genuine or has been tampered with (binary detection). At the \emph{pixel level},
it must also identify and delineate which specific regions have been altered
(forgery localization). Both objectives are operationally necessary.

Binary detection alone is insufficient for real-world forensic deployments.
A positive detection decision provides no \emph{actionable information} to the
investigator: it does not reveal whether the fraudster altered the portrait
photograph, falsified textual fields, or combined multiple manipulations. Without
spatial evidence, a human expert cannot determine the nature of the fraud or
build a legal case. Detection systems are also increasingly subject to
\emph{explainability requirements} under emerging AI regulations: a binary
``tampered'' label cannot satisfy the obligation
to provide a meaningful explanation of the automated decision. Conversely,
localization alone does not yield a reliable document-level confidence score
suitable for automated accept/reject decisions. A practical forensic system must
address both tasks simultaneously.

The primary technical difficulty in joint detection and localization is that the
two objectives are in \emph{structural tension}: the classification head
benefits from global image features, while the segmentation head requires
spatially precise, pixel-level representations. Prior work either treats the
two tasks independently \citep{trufor2023} or collapses both losses into a
scalar objective via a fixed weight $\lambda$~\citep{kendall2018multi}. The choice of $\lambda$ is arbitrary and dataset-dependent.
A second challenge is the risk of over-fitting the model-selection criterion:
hyperparameter optimization (HPO) that maximizes PR-AUC may sacrifice Dice, and vice versa.

This paper introduces DocVerify, a multi-task forensic document
analysis system that addresses both challenges. Our contributions are:

\begin{enumerate}
  \item \textit{Multi-objective Pareto HPO for multitask forensics.}
    Both PR-AUC and Dice are treated as simultaneous optimization objectives.
    The optimal configuration is selected from the Pareto front by minimizing
    the Euclidean distance to the ideal point~$(1,~1)$, requiring no manual
    task-weight tuning.

  \item \textit{DocVerify architecture.}
    An end-to-end multi-task CNN comprising a lightweight convolutional encoder
    coupled to a U-Net decoder for pixel-level mask prediction and a
    classification head for document-level detection, trained jointly with a
    combined BCE+Dice segmentation loss.

  \item \textit{Rigorous nested cross-validation.}
    A 10-fold outer $\times$ 5-fold inner nested cross-validation scheme with
    50 Multiobjective Tree-structured Parzen Estimator (MOTPE) trials per
    outer fold (2\,500 HPO evaluations total). Final performance is estimated
    on a blind holdout using 30 independent seeds, with Wilcoxon signed-rank
    tests~\citep{wilcoxon1945} and Holm--Bonferroni correction~\citep{holm1979}.
\end{enumerate}

The source code and experimental artifacts are publicly available at
\url{https://github.com/HernanDiaz/IDVerify/}.

\begin{figure*}[!t]
  \centering
  \includegraphics[width=\textwidth]{fig0_architecture_compact}
  \caption{DocVerify architecture. The encoder maps a $3{\times}224{\times}224$
    input to a $256{\times}7{\times}7$ bottleneck through six convolutional
    blocks ($8\!\to\!256$ channels). The U-Net decoder reconstructs the forgery
    mask via skip connections; the classification head outputs a document-level
    tampered/genuine score.}
  \label{fig:architecture}
\end{figure*}

The remainder of the paper is organized as follows. Section~\ref{sec:related}
reviews related work. Section~\ref{sec:method} describes the DocVerify
architecture and HPO framework. Section~\ref{sec:experiments} presents
experimental results. Section~\ref{sec:discussion} discusses the findings and
limitations. Section~\ref{sec:conclusion} concludes.

% ============================================================
% Section 2: Related Work
% ============================================================

\section{Related Work}
\label{sec:related}

\subsection{Document and Image Forgery Detection}

Deep learning has substantially advanced general-purpose image forensics:
CNNs can detect subtle statistical fingerprints left by editing operations,
often outperforming hand-crafted pipelines~\citep{mantranet2019}.
\citet{cozzolino2020noiseprint} (Noiseprint) train a CNN via Siamese
contrastive learning to extract camera-model noise residuals; forgeries appear
as anomalies in this residual without requiring semantic understanding.
\citet{trufor2023} (\textsc{TruFor}) generalize this with a transformer fusion
network (NoisePrint++) that jointly models RGB semantics and noise-pattern
anomalies, achieving strong performance on general image forensics benchmarks
with ${\sim}828\,000$ pre-training images. \revB{Closely related, deepfake and
face-manipulation detection has advanced rapidly in this journal:
\citet{liang2023twostream} pair a two-stream network with hierarchical
supervision, \citet{leporoni2024guided} fuse {RGB} and depth cues, and
\citet{wang2024deepfake} add a self-attention discriminator, each improving
robustness to synthetic facial forgeries.}

Identity document forgery introduces unique challenges: forged regions are
often small (a single text field or portrait), the structured background
(guilloché patterns, UV-reactive inks) creates strong false negatives for
general-purpose noise detectors, and document type diversity demands
generalization beyond natural-image pre-training. Official evaluations on
FantasyID confirm that general-purpose baselines (TruFor, MMFusion~\citep{mmfusion2024},
FatFormer~\citep{fatformer2024}, UniFD~\citep{unifd2023}) all achieve
false-negative rates above 50\% at a 10\% false-positive
threshold~\citep{fantasyid2024}, motivating domain-specific approaches.

Among systems trained exclusively on FantasyID, \citet{edgedoc2025} propose
EdgeDoc, a lightweight CNN-Transformer encoder (EdgeNeXt~\citep{edgenext2022}) fed with the green
channel and a NoisePrint residual, with a U-Net decoder optimized jointly for
classification and localization. \citet{remtkd2025} introduce Re-MTKD, a
reinforcement-learning multi-teacher knowledge distillation framework using a
ConvNeXt~v2 encoder; the Sunlight team adapted this to identity documents using
over 60\,000 additional external images, winning both tracks of the DeepID
2025 Challenge~\citep{deepid2025challenge}.

DocVerify differs from all prior document-forensics systems in three respects:
(i) it systematically searches the task-weight $\lambda_{\text{mask}}$ via
multi-objective HPO rather than fixing it heuristically; (ii) it selects the
final configuration via a Pareto-optimal criterion; and (iii) it reports
uncertainty through nested cross-validation, providing statistically rigorous
performance estimates.

\subsection{Multi-Task Learning and HPO}

In multi-task learning, the
dominant approach combines task losses via scalar scalarization~\citep{marler2004survey}
with fixed weights set before training. \citet{kendall2018multi} propose
learning these weights via homoscedastic uncertainty; \citet{liu2019end}
introduce gradient normalization. All scalar approaches commit to a fixed
trade-off before training begins.

Multi-objective HPO extends single-objective search to jointly optimize multiple
metrics~\citep{miettinen1999nonlinear}. The Multiobjective TPE
(MOTPE)~\citep{motpe2022} adapts the TPE surrogate~\citep{tpe2011} to the
multi-objective setting, achieving competitive Pareto-front coverage with far
fewer function evaluations than evolutionary methods such as
NSGA-II~\citep{nsga2_2002}, which is critical when each evaluation requires
training a deep network. To our knowledge, this is the first work to apply Pareto-based
multi-objective HPO to joint identity-document forgery detection and
localization within a nested cross-validation framework.

% ============================================================
% Section 3: Methodology
% ============================================================

\section{Methodology}
\label{sec:method}

\subsection{Dataset}
\label{subsec:dataset}

We train and evaluate DocVerify on FantasyID~\citep{fantasyid2024}, a public
benchmark for ID-document forgery detection and localization from Idiap for the
DeepID 2025 Challenge~\citep{deepid2025challenge}. The FantasyID release we
process contains 3\,284 labeled images of synthetic fantasy identity documents: \textit{bonafide}
(genuine) and \textit{attack} (tampered via face swap or text inpainting).

Each image is accompanied by JSON annotations specifying bounding-box regions
with \texttt{region\_provenance: altered}. We convert these to binary
segmentation masks at training resolution ($224 \times 224$ pixels). We
apply \texttt{GroupShuffleSplit} to all 3\,284 available FantasyID images
(grouping by document identity, $\text{test\_size}=0.15$,
\texttt{random\_state}=42) to obtain approximately 2\,791 development images
and 493 blind holdout images. \rev{This split was performed once,
before any hyperparameter search or model-design decision, and the holdout
was accessed only for the final blind evaluation
(Section~\ref{subsec:ablation}).} The dataset is class-imbalanced:
2\,351 attack images (71.6\%) and 933 bonafide images (28.4\%); this
imbalance is preserved in all splits and not corrected via resampling or
loss weighting. \rev{The mild imbalance (2.5:1, attack-majority)
leaves PR-AUC robust; a \texttt{pos\_weight} sweep confirms no significant
effect.}

\subsection{Architecture}
\label{subsec:arch}

DocVerify is a multi-task encoder-decoder network for simultaneous
document-level classification and pixel-level forgery localization
(Fig.~\ref{fig:architecture}).

\subsubsection{Encoder}
The encoder is a six-block lightweight CNN inspired by \citet{patel2023dcnn},
a deepfake-detection architecture specifically designed for forgery artifact
extraction. A lightweight encoder is a deliberate design choice: the nested
CV protocol requires training ${\sim}2\,500$ models, so per-trial cost must
be kept low. The channel progression is $8 \to 16 \to 32 \to 64 \to 128 \to
256$ filters. All layers use LeakyReLU ($\alpha=0.2$)
and batch normalization. Blocks 1--4 use $3 \times 3$
convolutions; blocks 5--6 use $5 \times 5$ convolutions. Downsampling uses
average pooling after blocks 2 and~3 (to preserve boundary information) and
max pooling after blocks 4 and~5 (for translation invariance).

\subsubsection{Classification Head}
Global average pooling is applied to the $7 \times 7$ bottleneck, yielding a
256-dimensional vector passed through four fully-connected layers
(256$\to$32$\to$16$\to$16$\to$1). \rev{Dropout (rate $p$) precedes each linear
layer and LeakyReLU ($\alpha=0.2$) each hidden layer; the final layer is linear,}
producing a classification logit $\hat{y}_{\text{cls}}$.

\subsubsection{Segmentation Head (U-Net Decoder)}
A five-stage U-Net decoder~\citep{unet2015} reconstructs the spatial mask.
\rev{Each stage applies bilinear $\times2$ upsampling, concatenates channel-wise
the matching-resolution encoder skip (tapped before pooling), then two $3\times3$
convolutions with LeakyReLU ($\alpha=0.2$) and no batch normalization.} The
channel width halves at each stage from
$C_{\text{dec}}$ (a tunable HPO parameter) down to $C_{\text{dec}}/16$. A final
$1 \times 1$ convolution produces the segmentation logit $\hat{Y}_{\text{mask}}$.

\subsection{Loss Functions}
\label{subsec:loss}

The combined training loss is:
\begin{equation}
  \mathcal{L} = \mathcal{L}_{\text{cls}} + \lambda_{\text{mask}} \,
                \mathcal{L}_{\text{seg}},
  \label{eq:total_loss}
\end{equation}
\rev{where, in Eq.~\ref{eq:total_loss},} $\lambda_{\text{mask}} \in [0.5, 3.0]$
is included in the HPO search
space. $\mathcal{L}_{\text{cls}} = \text{BCEWithLogits}(\hat{y}_{\text{cls}},
y_{\text{cls}})$. The segmentation loss combines BCE and Dice:
\begin{equation}
  \mathcal{L}_{\text{seg}} = \text{BCE}(\hat{Y}, Y) + 1 -
    \frac{2 \sum_i p_i g_i + \epsilon}{\sum_i p_i + \sum_i g_i + \epsilon},
  \label{eq:seg_loss}
\end{equation}
\rev{The Dice term in Eq.~\ref{eq:seg_loss},} as introduced for volumetric
segmentation~\citep{milletari2016vnet}, \rev{addresses}
foreground/background imbalance in small forgery regions.

All models are trained with AdamW and ReduceLROnPlateau
(factor $0.5$, patience~3). Mixed-precision training (FP16/FP32 via PyTorch AMP)
is used throughout.

\begin{figure*}[!t]
  \centering
  \includegraphics[width=\textwidth]{fig1_pareto_front}
  \caption{Validation objective space across all 500 HPO trials (10 folds
    $\times$ 50 MOTPE trials). Grey: evaluated; blue: Pareto-optimal;
    dashed: Pareto frontier. Red diamond: ideal point $(1,1)$; annotated point:
    selected configuration (distance $d=0.316$).}
  \label{fig:pareto}
\end{figure*}

\subsection{Multi-Objective HPO}
\label{subsec:hpo}

We optimize two complementary metrics: \emph{PR-AUC} (area under the
precision-recall curve, measuring document-level detection quality) and
\emph{Dice} ($2|P \cap G|/(|P|+|G|)$, measuring pixel-level mask overlap).
Our HPO formulation simultaneously maximizes:
\begin{equation}
  \max_{\boldsymbol{\theta} \in \Theta}\; \bigl(\text{PR-AUC}(\boldsymbol{\theta}),\;
  \text{Dice}(\boldsymbol{\theta})\bigr).
  \label{eq:moopt}
\end{equation}
\rev{In Eq.~\ref{eq:moopt},} $\Theta$ \rev{is the search space defined in
Table~\ref{tab:hpo_space}}.

\begin{table}[!t]
  \centering
  \caption{HPO search space \rev{and final selected configuration
           (Section~\ref{subsec:final_config}). All log-uniform except
           dropout and $\lambda_{\text{mask}}$ (uniform) and decoder
           channels}.}
  \label{tab:hpo_space}
  \begin{tabular}{llc}
    \toprule
    Parameter & Range & \rev{Selected} \\
    \midrule
    $\eta$ (learning rate) & $[5{\times}10^{-5},\; 9{\times}10^{-4}]$ & \rev{$5.21{\times}10^{-4}$} \\
    $\lambda_{\text{wd}}$ (weight decay) & $[10^{-7},\; 10^{-4}]$ & \rev{$4.27{\times}10^{-5}$} \\
    $p$ (dropout) & $[0.1,\; 0.4]$ & \rev{$0.103$} \\
    $C_{\text{dec}}$ (decoder channels) & $\{96, 128, 192, 256\}$ & \rev{$192$} \\
    $\lambda_{\text{mask}}$ (mask loss weight) & $[0.5,\; 3.0]$ & \rev{$2.46$} \\
    \bottomrule
  \end{tabular}
\end{table}

We use MOTPE~\citep{motpe2022} as implemented in Optuna~\citep{optuna2019}.
After 50 trials on the inner splits, we select the configuration
$\boldsymbol{\theta}^*$ from the Pareto front $\mathcal{P}$ by minimizing the
Euclidean distance to the ideal point \rev{(Eq.~\ref{eq:ideal_point})}:
\begin{equation}
  \boldsymbol{\theta}^* = \argmin_{\boldsymbol{\theta} \in \mathcal{P}}
    \sqrt{(1 - \text{PR-AUC})^2 + (1 - \text{Dice})^2}.
  \label{eq:ideal_point}
\end{equation}
This criterion requires no user-specified weighting and assigns equal importance
to detection and localization. Figure~\ref{fig:pareto} illustrates the
resulting objective space.

\subsection{Nested Cross-Validation}
\label{subsec:ncv}

To obtain unbiased estimates while keeping HPO blind to test data, we use a
three-level nested CV scheme \citep{varma2006bias}. The outer loop uses 10
stratified folds \rev{of the training set ($85\%$)}; each fold serves once as
the test set. For each outer fold,
the remaining data is split into 5 inner folds on which 50 MOTPE trials are
run (up to 15 epochs each, patience~$=1$). The optimal $\boldsymbol{\theta}^*$
is then used to retrain from scratch on all outer-fold training data (up to 100
epochs, early stopping patience~$=12$ \rev{on the validation distance to the
ideal point (Eq.~\ref{eq:ideal_point}), as in Pareto selection}). This design ensures HPO never sees test data: 50 MOTPE trials per outer fold
yield 500 unique hyperparameter configurations in total (10 outer
folds $\times$ 50 trials); because each configuration is evaluated on 5 inner
splits, this corresponds to $10 \times 50 \times 5 = 2\,500$ individual model
trainings, plus 10 final models retrained on full outer-fold data.

% ============================================================
% Section 4: Experiments and Results
% ============================================================

\section{Experiments and Results}
\label{sec:experiments}

\subsection{Experimental Setup}
\label{subsec:setup}

All experiments are conducted on a single workstation running Windows~11,
equipped with an AMD Ryzen~5 3400G (3.70\,GHz), 32\,GB DDR4-3200 RAM, and an NVIDIA
RTX~5060\,Ti (16\,GB VRAM), using PyTorch 2.10.0 with CUDA 13.0 and
Optuna~3.x. Training uses a batch size of 64 and no data augmentation
\rev{(standard augmentation, rotation, brightness, contrast, blur and JPEG,
significantly degraded every metric on the FantasyID holdout;
see~\citep{docverify_zenodo})}.
All images are resized to $224\!\times\!224$ pixels and normalised to
$[0,1]$ prior to caching in GPU memory.
The full NCV pipeline required
approximately 59~hours; VRAM caching reduces per-trial training time from
$\sim$12 to $\sim$3.5 minutes. \rev{This 59-hour budget covers the full
scientific protocol; retraining the single deployed model from scratch takes
$7.7 \pm 1.0$~min.}

\subsection{Nested Cross-Validation Results}
\label{subsec:ncv_results}

Table~\ref{tab:nested_cv} reports per-fold test metrics. PR-AUC is
consistently high across all folds ($0.977$--$0.999$, mean $0.9921 \pm
0.0058$). Dice shows slightly more fold-to-fold variability ($0.800$--$0.909$,
mean $0.856 \pm 0.030$), reflecting the difficulty of pixel-level localization.
The Pareto-selected $\lambda_{\text{mask}}$ clusters between 2.3 and 2.7 in
eight of ten folds; the two outlier folds (F5, F9) correspond to the two lowest
Dice scores, confirming that $\lambda_{\text{mask}}$ is a decisive driver of
localization quality.

\begin{table}[!t]
  \centering
  \caption{Per-fold test results from 10-fold NCV.
           $\overline{w}_{\text{mask}}$: selected loss weight.}
  \label{tab:nested_cv}
  \begin{tabular}{cccc}
    \toprule
    Fold & PR-AUC & Dice & $\overline{w}_{\text{mask}}$ \\
    \midrule
    F1  & 0.9939 & 0.861 & 2.46 \\
    F2  & 0.9769 & 0.871 & 2.70 \\
    F3  & 0.9993 & 0.909 & 2.47 \\
    F4  & 0.9915 & 0.858 & 2.43 \\
    F5  & 0.9948 & 0.827 & 1.41 \\
    F6  & 0.9936 & 0.874 & 2.89 \\
    F7  & 0.9940 & 0.861 & 2.67 \\
    F8  & 0.9967 & 0.887 & 2.29 \\
    F9  & 0.9896 & 0.807 & 1.95 \\
    F10 & 0.9909 & 0.800 & 2.65 \\
    \midrule
    \textbf{Mean (Std)} & \textbf{0.9921 (0.0058)} & \textbf{0.856 (0.030)} & \textbf{2.39 (0.43)} \\
    \bottomrule
  \end{tabular}
\end{table}

\begin{table*}[!t]
  \centering
  \caption{\rev{Ablation and selection-criterion comparison (blind test,
           30 seeds); all variants vs.\ \textbf{Pareto (ours)}. $p$: paired
           Wilcoxon~\citep{wilcoxon1945}\,+\,Holm~\citep{holm1979} (significant
           at $p<0.05$); $d$: Cohen's~$d$~\citep{cohen1988statistical}
           ($|d|\!\approx\!0.2$ small, $0.5$ medium, $\geq\!0.8$ large).}}
  \label{tab:ablation}
  \begin{tabular}{lcccccc}
    \toprule
    & \multicolumn{3}{c}{PR-AUC} & \multicolumn{3}{c}{Dice} \\
    \cmidrule(r){2-4}\cmidrule(l){5-7}
    Variant & Mean $\pm$ Std & $p$ & $d$ & Mean $\pm$ Std & $p$ & $d$ \\
    \midrule
    \textbf{Pareto (ours)} & $\mathbf{0.9967 \pm 0.0021}$ & — & — & $\mathbf{0.875 \pm 0.018}$ & — & — \\
    \rev{\texttt{by\_prauc}} & \rev{$0.9953 \pm 0.0021$} & \rev{$0.06$} & \rev{$0.49$} & \rev{$0.859 \pm 0.023$} & \rev{$0.04$} & \rev{$0.59$} \\
    \rev{\texttt{by\_dice}}  & \rev{$0.9948 \pm 0.0027$} & \rev{$0.007$} & \rev{$0.64$} & \rev{$0.857 \pm 0.015$} & \rev{$<0.001$} & \rev{$1.00$} \\
    \midrule
    \texttt{cls\_only}          & $0.9503 \pm 0.0744$ & $<0.001$ & $0.62$ & $0.002 \pm 0.008$ & $<0.001$ & $47.2$ \\
    \texttt{seg\_only}          & $0.7277 \pm 0.0528$ & $<0.001$ & $5.15$ & $0.867 \pm 0.023$ & $0.29$ & $0.32$ \\
    \texttt{unweighted\_losses} & $0.9958 \pm 0.0016$ & $0.08$ & $0.37$ & $0.857 \pm 0.016$ & $<0.001$ & $1.01$ \\
    \bottomrule
  \end{tabular}
\end{table*}

\subsection{Final Model Configuration}
\label{subsec:final_config}

\rev{The cross-validation above yields ten fold-specific configurations and
estimates generalization, but the deployed system is a single model. We select
the configuration with the smallest inner-CV distance to the ideal point
($1,1$) across all outer folds (Eq.~\ref{eq:ideal_point};
Table~\ref{tab:hpo_space}), which fixes both
the architecture (decoder width) and the Pareto-optimized loss weight
$\lambda_{\text{mask}}=2.46$, consistent with the $2.3$--$2.7$ cluster of
Section~\ref{subsec:ncv_results}. This final configuration is retrained from
scratch on the development set (up to 100 epochs with early stopping, as in
Section~\ref{subsec:ncv}) and evaluated on the held-out 15\% blind-test
set; we repeat the retraining over 30 independent seeds to quantify
run-to-run variability.}

\subsection{Ablation and Selection-Criterion Study}
\label{subsec:ablation}

To isolate each design decision, we evaluate this final model against five
variants under the same 30-seed blind-test protocol
(Table~\ref{tab:ablation}). Three variants probe the architecture and loss
weighting: \texttt{cls\_only} (classification head only); \texttt{seg\_only}
(segmentation head only, detection by thresholding mean mask probability);
and \texttt{unweighted\_losses} (full multitask, $\lambda_{\text{cls}} =
\lambda_{\text{mask}} = 0.5$, bypassing HPO). \rev{Two further variants probe
the HPO selection criterion applied to the \emph{same} 500 MOTPE trials:
\texttt{by\_prauc} (the trial maximizing validation PR-AUC) and
\texttt{by\_dice} (the trial maximizing validation Dice), against our Pareto
criterion (distance to the ideal point).}

Regarding the architecture and loss variants, three findings emerge.
First, joint multitask training is essential:
\texttt{cls\_only} produces near-zero Dice ($0.002 \pm 0.008$). Second, a
dedicated classification head is necessary: \texttt{seg\_only} achieves
competitive Dice ($0.867$) but PR-AUC drops to $0.728 \pm 0.053$, confirming
that the segmentation logit alone is an unreliable detection signal. Third,
Pareto-optimized weights outperform equal weights specifically on localization:
\texttt{unweighted\_losses} matches the full model on PR-AUC ($p = 0.083$) but
shows significantly lower Dice ($p = 1.3 \times 10^{-4}$, $d = 1.01$, large
effect), supporting the HPO contribution.

A key question for any multi-objective HPO framework is whether the added
complexity over simpler scalar alternatives yields measurable gains. \revB{The
advantage over scalar selection is real but partial and metric-specific rather
than uniform. Against \texttt{by\_dice}, Pareto is significantly better after Holm
correction on four of five metrics: PR-AUC ($+0.0019$, $d=0.64$, $p=0.007$), Dice
($+0.018$, $d=1.00$, $p<0.001$), balanced accuracy ($d=0.45$, $p=0.040$) and
F1-macro ($d=0.47$, $p=0.043$), though only the Dice gain reaches a large effect.
Against \texttt{by\_prauc} the advantage is narrower: only Dice ($d=0.59$,
$p=0.043$) and F1-macro ($d=0.46$, $p=0.050$, at the significance threshold) are
significant, while PR-AUC, balanced accuracy and F1-attack are not ($p>0.06$).
Effect sizes are otherwise small-to-medium ($d\approx0.4$--$0.65$) and the
absolute gains modest ($\Delta\lesssim0.02$), so the practical benefit is
concentrated in localization quality; its broader value is that distance-to-ideal
selection attains this without committing a priori to a primary objective (full
per-metric statistics in Supplementary Table~S4).}

Beyond predictive accuracy, Pareto selection requires no a priori commitment to a
primary objective, making the full trade-off landscape observable
(Fig.~\ref{fig:pareto}) and remaining structurally consistent with the DeepID
challenge protocol, which ranks detection and localization independently.

\begin{figure*}[!t]
  \centering
  \includegraphics[width=\textwidth]{fig_qualitative}
  \caption{\rev{Representative blind-test predictions (single seed). Top row:
  ground-truth altered regions (green); bottom row: regions predicted by
  DocVerify (red). Left: a forged photograph correctly detected and localized.
  Centre: a forgery correctly flagged but poorly localized. Right: a bonafide
  document misclassified as an attack.}}
  \label{fig:qualitative}
\end{figure*}

\subsection{Comparison with Related Systems}
\label{subsec:challenge}

\rev{DocVerify was developed after the DeepID 2025
Challenge~\citep{deepid2025challenge} deadline; we report results on our
internal 15\% holdout (30 seeds), which differs from the official challenge
set, so any comparison is indicative only. Training data also differ markedly. The leading
entries rely on external data: Sunlight (1st) augmented FantasyID with over
60\,000 identity images, while the third-place EdgeDoc+TruFor fusion and the
TruFor baseline~\citep{trufor2023}\footnote{\revB{For completeness, per-seed
scores for a TruFor model fine-tuned on FantasyID are released with the code
(\texttt{trufor\_finetuned\_scores.csv}); we omit them from the comparison here
because, like the challenge figures, they were not obtained on our internal
holdout, so no controlled comparison is possible.}} build on a backbone
pre-trained on $\sim$828\,000 manipulation images. EdgeDoc standalone~\citep{edgedoc2025},
UAM-Biometrics, and DocVerify instead train on FantasyID alone. With these
caveats, DocVerify reaches
detection F1@0.5 \rev{(the F1 score, $2PR/(P{+}R)$, of the bonafide/attack
decision at classification threshold $0.5$)} $=0.969 \pm 0.014$ and localization
F1$_{\text{loc}}$ \rev{(pixel-level F1 of the predicted mask)}
$= 0.807 \pm 0.096$. Reported challenge figures span EdgeDoc standalone
($0.43$~\citep{edgedoc2025}), UAM-Biometrics ($0.712$), the EdgeDoc+TruFor fusion
($0.958$~\citep{deepid2025challenge}, on the pre-trained TruFor backbone) and the
winner Sunlight ($0.991$, F1$_{\text{loc}}$ $0.784$). As evaluation sets and
training data differ, we make no superiority claim and imply no ranking; the
figures only add context.}

\subsection{\rev{Generalization}}
\label{subsec:generalization}

\rev{Two experiments probe whether the approach generalizes beyond its specific
setting. \emph{Across architectures}, we re-ran the identical MOTPE protocol
(50 trials, 5 inner splits, Pareto selection, 30-seed blind test; a single study
rather than the full nested loop, to bound cost) with the encoder replaced by an
ImageNet-pretrained ResNet-18 feeding the same U-Net decoder. MOTPE selected
$\lambda_{\text{mask}} = 2.69$ (vs.\ $2.46$ for DocVerify), so the Pareto balance
is consistent across backbones. It reaches PR-AUC $0.994 \pm 0.003$ and Dice
$0.889 \pm 0.011$ (30 seeds); paired Wilcoxon\,+\,Holm gives DocVerify a small
PR-AUC edge ($d = 0.75$) and ResNet-18 a Dice/F1 edge ($d = 0.60$--$0.67$), with
balanced accuracy tied. Neither dominates, indicating the selection methodology
is not tied to a specific backbone~\citep{docverify_zenodo}. \emph{Across datasets},
we evaluate the FantasyID blind-test models zero-shot on the out-of-domain SIDTD
dataset~\citep{sidtd2024}, using only its \emph{templates} subset (2\,222 images:
1\,000 bonafide, 1\,222 attack) rather than the full SIDTD video collection.
Zero-shot performance drops to near chance (PR-AUC $0.555 \pm 0.011$, ROC-AUC
$0.504 \pm 0.009$): the FantasyID$\rightarrow$SIDTD domain gap is too wide for a
model trained on FantasyID alone. Retraining the same architecture on an 80/20
stratified split of the templates subset (1\,777 train, 445 test) recovers PR-AUC
$0.878 \pm 0.016$, ROC-AUC $0.826 \pm 0.023$, whereas a few-shot sweep
(10--400 images) stays at chance: only full in-domain training closes the gap, so
it reflects domain shift, not architectural capacity. The zero-shot scope of the
current FantasyID-trained system is thus bounded to its training domain.}

\subsection{Qualitative Analysis}
\label{subsec:qualitative}

\rev{Figure~\ref{fig:qualitative} contrasts three blind-test outcomes: a forged
photograph precisely localized (left, $p_{\text{attack}}=0.96$, Dice $0.96$), and
the two failure modes, a forgery correctly flagged but with its mask drifting onto
the portrait (centre, $p_{\text{attack}}=1.00$, Dice $0.22$) and a bonafide
document wrongly flagged with a spurious mask (right, $p_{\text{attack}}=0.96$),
illustrating the model's localization variance and false positives.}

% ============================================================
% Section 5: Discussion
% ============================================================

\section{Discussion}
\label{sec:discussion}

\subsection{Analysis of Key Findings}

The ablation reveals a task asymmetry: \texttt{seg\_only} achieves competitive
Dice ($0.867 \pm 0.023$, $p = 0.288$ vs.\ Pareto) but its PR-AUC collapses to
$0.728 \pm 0.053$. The U-Net decoder captures forgery evidence for pixel-level
prediction yet cannot aggregate it into a reliable document-level decision,
confirming the shared encoder, regularized by the segmentation loss, learns
spatially-aware features that improve detection.

Comparing Pareto against \texttt{unweighted\_losses} isolates the HPO
loss-weighting contribution: these differ only in whether $\lambda_{\text{mask}}$
is HPO-set \rev{($2.46$ deployed, NCV fold mean $2.39 \pm 0.43$)} or fixed at
$0.5$. Pareto-optimized weights yield significantly higher Dice ($p =
1.3 \times 10^{-4}$, $d = 1.01$) at equal PR-AUC, confirming
$\lambda_{\text{mask}} \approx 2.4$ is broadly optimal. \rev{The comparison
against the two scalar criteria (Sec.~\ref{subsec:ablation}) further shows the
gain stems from the selection rule itself: scalar criteria over-commit to one
axis, whereas distance-to-ideal selection balances both objectives without a
priori weighting.}

\subsection{Limitations}

\paragraph{\rev{Computational cost.}}
\rev{The full nested-CV protocol (2{,}500 trainings, ${\sim}59$~h) is computationally
heavy, but this reflects the unbiased-estimation protocol, not MOTPE itself:
HPO alone can be run on a single train/validation split at a fraction of the
cost, and the deployed model retrains in $7.7 \pm 1.0$~min.}

\paragraph{\rev{Generalization.}}
\rev{Trained on FantasyID alone, DocVerify falls to near chance zero-shot on SIDTD
(Sec.~\ref{subsec:generalization}); even the top challenge entries generalized
only by training on large external corpora (Sec.~\ref{subsec:challenge}), making
multi-source training the main future direction.} \revB{Standard augmentation does
not close this gap: augmentation-trained models remain at near chance zero-shot on
SIDTD (Supplementary Table~S8), indicating the gap reflects domain shift rather
than limited input variability.}

% ============================================================
% Section 6: Conclusion
% ============================================================

\section{Conclusion}
\label{sec:conclusion}

We presented DocVerify, a multi-task CNN for joint detection and localization of
forged identity documents, paired with a multi-objective Pareto HPO framework
that selects hyperparameter configurations by optimizing PR-AUC and Dice
simultaneously without fixing their relative importance in advance.

Our main findings are: \rev{(i) learning detection and localization jointly
benefits both tasks: a single-task model performs well only on the axis it
targets (detection-only Dice $0.002$; localization-only PR-AUC $0.728$; both
$p < 0.001$), whereas the joint model attains both; it further yields higher
Dice than equal-weight training ($p = 1.3 \times 10^{-4}$, $d = 1.01$) at no
PR-AUC cost,} \revB{and distance-to-ideal selection improves on scalar PR-AUC- or
Dice-maximizing selection chiefly in localization, a partial and metric-specific
edge}; (ii) under a rigorous 10-fold
nested cross-validation protocol, DocVerify achieves PR-AUC $= 0.9921 \pm
0.0058$ and Dice $= 0.856 \pm 0.030$, demonstrating consistent generalization
across data partitions; and \rev{(iii) trained on FantasyID alone, DocVerify
attains F1@0.5 $= 0.969 \pm 0.014$ and F1$_{\text{loc}} = 0.807 \pm 0.096$ on our
internal holdout, reported alongside the DeepID 2025 Challenge for context only,
as the evaluation sets differ.}

\revB{Two strengths and two limitations frame these results. The multi-objective
formulation removes the need to hand-tune the detection/localization trade-off, and
the same MOTPE protocol transfers to a ResNet-18 backbone, so the contribution is
methodological rather than architecture-specific; the evaluation is fully
reproducible, with per-seed results and extended tables archived openly. Set against
this, the benefit of Pareto selection over simple scalar criteria is partial and
concentrated in localization rather than uniform across metrics, and the reported
accuracy is bounded to the in-domain FantasyID setting: zero-shot transfer to SIDTD
is near chance and is not recovered by standard augmentation, so cross-domain
deployment would require multi-source training. The nested protocol is also
computationally heavy, though a single deployed model retrains in minutes.}

\rev{Applying the identical MOTPE protocol to a ResNet-18 backbone yields
comparable blind-test performance (Sec.~\ref{subsec:generalization}), so the
approach is not tied to one architecture. By contrast, zero-shot evaluation on
SIDTD exposed a substantial cross-domain gap; closing it through multi-source
training is a priority for future work}, alongside forgery-artifact features
(noise fingerprints, frequency anomalies~\citep{cozzolino2020noiseprint}) and
foundation model backbones~\citep{fakeidet2025} for richer pre-training.

\section*{Acknowledgments}
This work was supported by the Spanish National Cybersecurity Institute (INCIBE,
Spain). The author thanks the Idiap Research Institute for releasing the
FantasyID dataset.

% ============================================================
% BIBLIOGRAPHY
% ============================================================

%% Bibliography inlined from .bbl (self-contained; no bibtex/.bib needed)
\begin{thebibliography}{32}
\expandafter\ifx\csname natexlab\endcsname\relax\def\natexlab#1{#1}\fi
\providecommand{\url}[1]{\texttt{#1}}
\providecommand{\href}[2]{#2}
\providecommand{\path}[1]{#1}
\providecommand{\DOIprefix}{doi:}
\providecommand{\ArXivprefix}{arXiv:}
\providecommand{\URLprefix}{URL: }
\providecommand{\Pubmedprefix}{pmid:}
\providecommand{\doi}[1]{\href{http://dx.doi.org/#1}{\path{#1}}}
\providecommand{\Pubmed}[1]{\href{pmid:#1}{\path{#1}}}
\providecommand{\bibinfo}[2]{#2}
\ifx\xfnm\relax \def\xfnm[#1]{\unskip,\space#1}\fi
%Type = Inproceedings
\bibitem[{Guillaro et~al.(2023)Guillaro, Cozzolino, Sud, Dufour, and
  Verdoliva}]{trufor2023}
\bibinfo{author}{F.~Guillaro}, \bibinfo{author}{D.~Cozzolino},
  \bibinfo{author}{A.~Sud}, \bibinfo{author}{N.~Dufour},
  \bibinfo{author}{L.~Verdoliva},
\newblock \bibinfo{title}{{TruFor}: Leveraging all-round clues for trustworthy
  image forgery detection and localization},
\newblock in: \bibinfo{booktitle}{Proceedings of the {IEEE/CVF} Conference on
  Computer Vision and Pattern Recognition (CVPR)}, \bibinfo{year}{2023}, pp.
  \bibinfo{pages}{20606--20615}.
%Type = Inproceedings
\bibitem[{Cipolla et~al.(2018)Cipolla, Gal, and Kendall}]{kendall2018multi}
\bibinfo{author}{R.~Cipolla}, \bibinfo{author}{Y.~Gal},
  \bibinfo{author}{A.~Kendall},
\newblock \bibinfo{title}{Multi-task learning using uncertainty to weigh losses
  for scene geometry and semantics},
\newblock in: \bibinfo{booktitle}{2018 IEEE/CVF Conference on Computer Vision
  and Pattern Recognition}, \bibinfo{publisher}{IEEE}, \bibinfo{year}{2018},
  pp. \bibinfo{pages}{7482--7491}. \DOIprefix\doi{10.1109/cvpr.2018.00781}.
%Type = Article
\bibitem[{Wilcoxon(1945)}]{wilcoxon1945}
\bibinfo{author}{F.~Wilcoxon},
\newblock \bibinfo{title}{Individual comparisons by ranking methods},
\newblock \bibinfo{journal}{Biometrics Bulletin} \bibinfo{volume}{1}
  (\bibinfo{year}{1945}) \bibinfo{pages}{80--83}.
  \DOIprefix\doi{10.2307/3001968}.
%Type = Article
\bibitem[{Holm(1979)}]{holm1979}
\bibinfo{author}{S.~Holm},
\newblock \bibinfo{title}{A simple sequentially rejective multiple test
  procedure},
\newblock \bibinfo{journal}{Scandinavian Journal of Statistics}
  \bibinfo{volume}{6} (\bibinfo{year}{1979}) \bibinfo{pages}{65--70}.
%Type = Inproceedings
\bibitem[{Wu et~al.(2019)Wu, AbdAlmageed, and Natarajan}]{mantranet2019}
\bibinfo{author}{Y.~Wu}, \bibinfo{author}{W.~AbdAlmageed},
  \bibinfo{author}{P.~Natarajan},
\newblock \bibinfo{title}{Mantra-net: Manipulation tracing network for
  detection and localization of image forgeries with anomalous features},
\newblock in: \bibinfo{booktitle}{2019 IEEE/CVF Conference on Computer Vision
  and Pattern Recognition (CVPR)}, \bibinfo{publisher}{IEEE},
  \bibinfo{year}{2019}, pp. \bibinfo{pages}{9535--9544}.
  \DOIprefix\doi{10.1109/cvpr.2019.00977}.
%Type = Article
\bibitem[{Cozzolino and Verdoliva(2020)}]{cozzolino2020noiseprint}
\bibinfo{author}{D.~Cozzolino}, \bibinfo{author}{L.~Verdoliva},
\newblock \bibinfo{title}{Noiseprint: A cnn-based camera model fingerprint},
\newblock \bibinfo{journal}{IEEE Transactions on Information Forensics and
  Security} \bibinfo{volume}{15} (\bibinfo{year}{2020})
  \bibinfo{pages}{144--159}. \DOIprefix\doi{10.1109/tifs.2019.2916364}.
%Type = Article
\bibitem[{Liang et~al.(2023)Liang, Wang, Jin, Pan, and
  Liu}]{liang2023twostream}
\bibinfo{author}{Y.~Liang}, \bibinfo{author}{M.~Wang},
  \bibinfo{author}{Y.~Jin}, \bibinfo{author}{S.~Pan}, \bibinfo{author}{Y.~Liu},
\newblock \bibinfo{title}{Hierarchical supervisions with two-stream network for
  {Deepfake} detection},
\newblock \bibinfo{journal}{Pattern Recognition Letters} \bibinfo{volume}{172}
  (\bibinfo{year}{2023}) \bibinfo{pages}{121--127}.
  \DOIprefix\doi{10.1016/j.patrec.2023.05.029}.
%Type = Article
\bibitem[{Leporoni et~al.(2024)Leporoni, Maiano, Papa, and
  Amerini}]{leporoni2024guided}
\bibinfo{author}{G.~Leporoni}, \bibinfo{author}{L.~Maiano},
  \bibinfo{author}{L.~Papa}, \bibinfo{author}{I.~Amerini},
\newblock \bibinfo{title}{A guided-based approach for deepfake detection:
  {RGB}-depth integration via features fusion},
\newblock \bibinfo{journal}{Pattern Recognition Letters} \bibinfo{volume}{181}
  (\bibinfo{year}{2024}) \bibinfo{pages}{99--105}.
  \DOIprefix\doi{10.1016/j.patrec.2024.03.025}.
%Type = Article
\bibitem[{Wang et~al.(2024)Wang, Zhu, Chen, Bi, and Wang}]{wang2024deepfake}
\bibinfo{author}{S.~Wang}, \bibinfo{author}{D.~Zhu}, \bibinfo{author}{J.~Chen},
  \bibinfo{author}{J.~Bi}, \bibinfo{author}{W.~Wang},
\newblock \bibinfo{title}{Deepfake face discrimination based on self-attention
  mechanism},
\newblock \bibinfo{journal}{Pattern Recognition Letters} \bibinfo{volume}{183}
  (\bibinfo{year}{2024}) \bibinfo{pages}{92--97}.
  \DOIprefix\doi{10.1016/j.patrec.2024.02.019}.
%Type = Inproceedings
\bibitem[{Triaridis et~al.(2024)Triaridis, Tsigos, and Mezaris}]{mmfusion2024}
\bibinfo{author}{K.~Triaridis}, \bibinfo{author}{K.~Tsigos},
  \bibinfo{author}{V.~Mezaris},
\newblock \bibinfo{title}{{MMFusion}: Combining image forensic filters for
  visual manipulation detection and localization},
\newblock in: \bibinfo{booktitle}{Proceedings of the 30th International
  Conference on MultiMedia Modeling (MMM)}, \bibinfo{year}{2024}.
%Type = Inproceedings
\bibitem[{Liu et~al.(2024)Liu, Tan, Tan, Wei, Zhao, and Wang}]{fatformer2024}
\bibinfo{author}{H.~Liu}, \bibinfo{author}{Z.~Tan}, \bibinfo{author}{C.~Tan},
  \bibinfo{author}{Y.~Wei}, \bibinfo{author}{Y.~Zhao},
  \bibinfo{author}{J.~Wang},
\newblock \bibinfo{title}{Forgery-aware adaptive transformer for generalizable
  synthetic image detection},
\newblock in: \bibinfo{booktitle}{Proceedings of the {IEEE/CVF} Conference on
  Computer Vision and Pattern Recognition (CVPR)}, \bibinfo{year}{2024}, pp.
  \bibinfo{pages}{10770--10780}.
%Type = Inproceedings
\bibitem[{Ojha et~al.(2023)Ojha, Li, and Lee}]{unifd2023}
\bibinfo{author}{U.~Ojha}, \bibinfo{author}{Y.~Li}, \bibinfo{author}{Y.~J.
  Lee},
\newblock \bibinfo{title}{Towards universal fake image detectors that
  generalize across generative models},
\newblock in: \bibinfo{booktitle}{Proceedings of the {IEEE/CVF} Conference on
  Computer Vision and Pattern Recognition (CVPR)}, \bibinfo{year}{2023}, pp.
  \bibinfo{pages}{24480--24489}.
%Type = Inproceedings
\bibitem[{Korshunov et~al.(2025)Korshunov, Mohammadi, Vidit, Ecabert, and
  Marcel}]{fantasyid2024}
\bibinfo{author}{P.~Korshunov}, \bibinfo{author}{A.~Mohammadi},
  \bibinfo{author}{V.~Vidit}, \bibinfo{author}{C.~Ecabert},
  \bibinfo{author}{S.~Marcel},
\newblock \bibinfo{title}{{FantasyID}: A dataset for detecting digital
  manipulations of {ID}-documents},
\newblock in: \bibinfo{booktitle}{Proceedings of the IEEE International Joint
  Conference on Biometrics (IJCB)}, \bibinfo{year}{2025}.
%Type = Misc
\bibitem[{George and Marcel(2025)}]{edgedoc2025}
\bibinfo{author}{A.~George}, \bibinfo{author}{S.~Marcel},
  \bibinfo{title}{{EdgeDoc}: Hybrid {CNN}-transformer model for accurate
  forgery detection and localization in {ID} documents}, \bibinfo{year}{2025}.
  \href{http://arxiv.org/abs/2508.16284}{{\tt arXiv:2508.16284}},
  \bibinfo{note}{idiap Research Report RR-08-2025; 3rd place DeepID 2025
  Challenge}.
%Type = Inproceedings
\bibitem[{Maaz et~al.(2022)Maaz, Shaker, Cholakkal, Khan, Zamir, Anwer, and
  Khan}]{edgenext2022}
\bibinfo{author}{M.~Maaz}, \bibinfo{author}{A.~Shaker},
  \bibinfo{author}{H.~Cholakkal}, \bibinfo{author}{S.~Khan},
  \bibinfo{author}{S.~W. Zamir}, \bibinfo{author}{R.~M. Anwer},
  \bibinfo{author}{F.~S. Khan},
\newblock \bibinfo{title}{{EdgeNeXt}: Efficiently amalgamated {CNN}-transformer
  architecture for mobile vision applications},
\newblock in: \bibinfo{booktitle}{European Conference on Computer Vision
  Workshops (ECCVW)}, \bibinfo{year}{2022}, pp. \bibinfo{pages}{3--20}.
%Type = Article
\bibitem[{Yu et~al.(2025)Yu, Ni, Zhang, Deng, and Lin}]{remtkd2025}
\bibinfo{author}{Z.~Yu}, \bibinfo{author}{J.~Ni}, \bibinfo{author}{J.~Zhang},
  \bibinfo{author}{H.~Deng}, \bibinfo{author}{Y.~Lin},
\newblock \bibinfo{title}{Reinforced multi-teacher knowledge distillation for
  efficient general image forgery detection and localization},
\newblock \bibinfo{journal}{Proceedings of the AAAI Conference on Artificial
  Intelligence} \bibinfo{volume}{39} (\bibinfo{year}{2025})
  \bibinfo{pages}{995--1003}. \DOIprefix\doi{10.1609/aaai.v39i1.32085}.
%Type = Inproceedings
\bibitem[{Korshunov et~al.(2025)Korshunov, Vidit, Mohammadi, Ecabert, Shamoska,
  Marcel et~al.}]{deepid2025challenge}
\bibinfo{author}{P.~Korshunov}, \bibinfo{author}{V.~Vidit},
  \bibinfo{author}{A.~Mohammadi}, \bibinfo{author}{C.~Ecabert},
  \bibinfo{author}{N.~Shamoska}, \bibinfo{author}{S.~Marcel}, et~al.,
\newblock \bibinfo{title}{{DeepID}: Challenge of detecting synthetic
  manipulations in {ID} documents},
\newblock in: \bibinfo{booktitle}{Proceedings of the {IEEE/CVF} International
  Conference on Computer Vision Workshops ({ICCVW})}, \bibinfo{year}{2025}, pp.
  \bibinfo{pages}{510--519}.
%Type = Article
\bibitem[{Marler and Arora(2004)}]{marler2004survey}
\bibinfo{author}{R.~T. Marler}, \bibinfo{author}{J.~S. Arora},
\newblock \bibinfo{title}{Survey of multi-objective optimization methods for
  engineering},
\newblock \bibinfo{journal}{Structural and Multidisciplinary Optimization}
  \bibinfo{volume}{26} (\bibinfo{year}{2004}) \bibinfo{pages}{369--395}.
  \DOIprefix\doi{10.1007/s00158-003-0368-6}.
%Type = Inproceedings
\bibitem[{Liu et~al.(2019)Liu, Johns, and Davison}]{liu2019end}
\bibinfo{author}{S.~Liu}, \bibinfo{author}{E.~Johns}, \bibinfo{author}{A.~J.
  Davison},
\newblock \bibinfo{title}{End-to-end multi-task learning with attention},
\newblock in: \bibinfo{booktitle}{2019 IEEE/CVF Conference on Computer Vision
  and Pattern Recognition (CVPR)}, \bibinfo{publisher}{IEEE},
  \bibinfo{year}{2019}, pp. \bibinfo{pages}{1871--1880}.
  \DOIprefix\doi{10.1109/cvpr.2019.00197}.
%Type = Book
\bibitem[{Miettinen(1999)}]{miettinen1999nonlinear}
\bibinfo{author}{K.~Miettinen}, \bibinfo{title}{Nonlinear Multiobjective
  Optimization}, \bibinfo{publisher}{Springer US}, \bibinfo{year}{1999}.
  \DOIprefix\doi{10.1007/978-1-4615-5563-6}.
%Type = Article
\bibitem[{Ozaki et~al.(2022)Ozaki, Tanigaki, Watanabe, Nomura, and
  Onishi}]{motpe2022}
\bibinfo{author}{Y.~Ozaki}, \bibinfo{author}{Y.~Tanigaki},
  \bibinfo{author}{S.~Watanabe}, \bibinfo{author}{M.~Nomura},
  \bibinfo{author}{M.~Onishi},
\newblock \bibinfo{title}{Multiobjective tree-structured parzen estimator},
\newblock \bibinfo{journal}{Journal of Artificial Intelligence Research}
  \bibinfo{volume}{73} (\bibinfo{year}{2022}) \bibinfo{pages}{1209--1250}.
  \DOIprefix\doi{10.1613/jair.1.13188}.
%Type = Inproceedings
\bibitem[{Bergstra et~al.(2011)Bergstra, Bardenet, Bengio, and
  K{\'e}gl}]{tpe2011}
\bibinfo{author}{J.~Bergstra}, \bibinfo{author}{R.~Bardenet},
  \bibinfo{author}{Y.~Bengio}, \bibinfo{author}{B.~K{\'e}gl},
\newblock \bibinfo{title}{Algorithms for hyper-parameter optimization},
\newblock in: \bibinfo{booktitle}{Advances in Neural Information Processing
  Systems (NeurIPS)}, volume~\bibinfo{volume}{24}, \bibinfo{year}{2011}.
%Type = Article
\bibitem[{Deb et~al.(2002)Deb, Pratap, Agarwal, and Meyarivan}]{nsga2_2002}
\bibinfo{author}{K.~Deb}, \bibinfo{author}{A.~Pratap},
  \bibinfo{author}{S.~Agarwal}, \bibinfo{author}{T.~Meyarivan},
\newblock \bibinfo{title}{A fast and elitist multiobjective genetic algorithm:
  Nsga-ii},
\newblock \bibinfo{journal}{IEEE Transactions on Evolutionary Computation}
  \bibinfo{volume}{6} (\bibinfo{year}{2002}) \bibinfo{pages}{182--197}.
  \DOIprefix\doi{10.1109/4235.996017}.
%Type = Article
\bibitem[{Patel et~al.(2023)Patel, Tanwar, Bhattacharya, Gupta, Alsuwian,
  Davidson, and Mazibuko}]{patel2023dcnn}
\bibinfo{author}{Y.~Patel}, \bibinfo{author}{S.~Tanwar},
  \bibinfo{author}{P.~Bhattacharya}, \bibinfo{author}{R.~Gupta},
  \bibinfo{author}{T.~Alsuwian}, \bibinfo{author}{I.~E. Davidson},
  \bibinfo{author}{T.~F. Mazibuko},
\newblock \bibinfo{title}{An improved dense cnn architecture for deepfake image
  detection},
\newblock \bibinfo{journal}{IEEE Access} \bibinfo{volume}{11}
  (\bibinfo{year}{2023}) \bibinfo{pages}{22081--22095}.
  \DOIprefix\doi{10.1109/access.2023.3251417}.
%Type = Inbook
\bibitem[{Ronneberger et~al.(2015)Ronneberger, Fischer, and Brox}]{unet2015}
\bibinfo{author}{O.~Ronneberger}, \bibinfo{author}{P.~Fischer},
  \bibinfo{author}{T.~Brox}, \bibinfo{title}{U-Net: Convolutional Networks for
  Biomedical Image Segmentation}, \bibinfo{publisher}{Springer International
  Publishing}, \bibinfo{year}{2015}, pp. \bibinfo{pages}{234--241}.
  \DOIprefix\doi{10.1007/978-3-319-24574-4_28}.
%Type = Inproceedings
\bibitem[{Milletari et~al.(2016)Milletari, Navab, and
  Ahmadi}]{milletari2016vnet}
\bibinfo{author}{F.~Milletari}, \bibinfo{author}{N.~Navab},
  \bibinfo{author}{S.-A. Ahmadi},
\newblock \bibinfo{title}{{V-Net}: Fully convolutional neural networks for
  volumetric medical image segmentation},
\newblock in: \bibinfo{booktitle}{Proceedings of the 4th International
  Conference on 3D Vision (3DV)}, \bibinfo{year}{2016}, pp.
  \bibinfo{pages}{565--571}. \DOIprefix\doi{10.1109/3DV.2016.79}.
%Type = Inproceedings
\bibitem[{Akiba et~al.(2019)Akiba, Sano, Yanase, Ohta, and Koyama}]{optuna2019}
\bibinfo{author}{T.~Akiba}, \bibinfo{author}{S.~Sano},
  \bibinfo{author}{T.~Yanase}, \bibinfo{author}{T.~Ohta},
  \bibinfo{author}{M.~Koyama},
\newblock \bibinfo{title}{Optuna: A next-generation hyperparameter optimization
  framework},
\newblock in: \bibinfo{booktitle}{Proceedings of the 25th ACM SIGKDD
  International Conference on Knowledge Discovery \& Data Mining}, KDD ’19,
  \bibinfo{publisher}{ACM}, \bibinfo{year}{2019}, pp.
  \bibinfo{pages}{2623--2631}. \DOIprefix\doi{10.1145/3292500.3330701}.
%Type = Article
\bibitem[{Varma and Simon(2006)}]{varma2006bias}
\bibinfo{author}{S.~Varma}, \bibinfo{author}{R.~Simon},
\newblock \bibinfo{title}{Bias in error estimation when using cross-validation
  for model selection},
\newblock \bibinfo{journal}{BMC Bioinformatics} \bibinfo{volume}{7}
  (\bibinfo{year}{2006}). \DOIprefix\doi{10.1186/1471-2105-7-91}.
%Type = Misc
\bibitem[{D\'{i}az~Rodr\'{i}guez(2026)}]{docverify_zenodo}
\bibinfo{author}{H.~D\'{i}az~Rodr\'{i}guez}, \bibinfo{title}{{DocVerify}:
  Reproducible per-seed results and extended tables for multi-objective {HPO}
  of forged-{ID} detection}, \bibinfo{year}{2026}.
  \DOIprefix\doi{10.5281/zenodo.21076490}.
%Type = Book
\bibitem[{Cohen(1988)}]{cohen1988statistical}
\bibinfo{author}{J.~Cohen}, \bibinfo{title}{Statistical Power Analysis for the
  Behavioral Sciences}, \bibinfo{edition}{2nd} ed.,
  \bibinfo{publisher}{Lawrence Erlbaum Associates}, \bibinfo{year}{1988}.
%Type = Article
\bibitem[{Boned et~al.(2024)Boned, Talarmain, Ghanmi, Chiron, Biswas, Awal, and
  Ramos~Terrades}]{sidtd2024}
\bibinfo{author}{C.~Boned}, \bibinfo{author}{M.~Talarmain},
  \bibinfo{author}{N.~Ghanmi}, \bibinfo{author}{G.~Chiron},
  \bibinfo{author}{S.~Biswas}, \bibinfo{author}{A.~M. Awal},
  \bibinfo{author}{O.~Ramos~Terrades},
\newblock \bibinfo{title}{Synthetic dataset of {ID} and travel documents},
\newblock \bibinfo{journal}{Scientific Data} \bibinfo{volume}{11}
  (\bibinfo{year}{2024}) \bibinfo{pages}{1356}.
  \DOIprefix\doi{10.1038/s41597-024-04160-9}.
%Type = Inproceedings
\bibitem[{Mu{\~n}oz-Haro et~al.(2025)Mu{\~n}oz-Haro, Tolosana, Vera-Rodriguez,
  Morales, and Fierrez}]{fakeidet2025}
\bibinfo{author}{J.~Mu{\~n}oz-Haro}, \bibinfo{author}{R.~Tolosana},
  \bibinfo{author}{R.~Vera-Rodriguez}, \bibinfo{author}{A.~Morales},
  \bibinfo{author}{J.~Fierrez},
\newblock \bibinfo{title}{{FakeIDet}: Exploring patches for privacy-preserving
  fake {ID} detection},
\newblock in: \bibinfo{booktitle}{Proceedings of the IEEE International Joint
  Conference on Biometrics (IJCB)}, \bibinfo{year}{2025}.

\end{thebibliography}

\end{document}
